\documentclass[10pt]{article}

\usepackage[utf8]{inputenc}

\usepackage[T1]{fontenc}

\usepackage{amsmath}

\usepackage{amsfonts}

\usepackage{amssymb}

\usepackage[version=4]{mhchem}

\usepackage{stmaryrd}

\usepackage{mathrsfs}

\usepackage{bbold}

\usepackage{graphicx}

\usepackage[export]{adjustbox}

\graphicspath{ {./images/} }



\begin{document}

согласующейся со структурой его аддитивной группы и инвариантной относительно операций умножения на ненулевые элементы из $K$, является Т. в. І. (этим условиям удовлетворяет, в частности, дискретная топология в $E$ ). Т. в. п. над полями с дискретной топологией наз. то п о л о г и ч е с к и ми в е к то рв ы м и г р у п п а м и. 4. Пусть $E$ - векторное пространство над т. п. $K, \mathscr{P}-$ нек-рое множество полунорм на $E$. III а р о м радиуса $r>0$ по полунорме $p$ на $E$ наз. множество $\{x \in E, p(x)<r\}$. Множество пересечений всевозможных конечных семейств шаров [всевозможных (положительных) радиусов] по (всевозможным) полунормам, принадлежащим множеству $\mathscr{P}$, образует базис окрестностей нуля нек-рой топологии $\tau_{\mathscr{P}}$ в $E$, согласующейся со структурой векторного пространства; говорят, что эта топология задается, или определяется семейством полунорм $\mathscr{P}$. Если $K=\mathbb{R}$ или $K=\mathbb{C}$, то топология $\tau_{\mathscr{P}}$ локально выпукла; обратно, топология всякого локально выпуклого пространства (л. в. п.) $E$ может быть задана нек-рым множеством полунорм - нашример, множеством калибровочных функций (функционалов Минковского) произвольной предбазы окрестностей нуля в $E$, состоящей из закругленных выпуклых множеств.



Подмножество Т. в. П. наз. о г р а ни ч е н н ы м, если оно поглощается всякой окрестностью нуля.



Т. в. п. наз. н о р м и р у е м ы м, если его топология может быть задана одной нормой. Для того чтобы T. в. п. над полем $\mathbb{R}$ или $\mathbb{C}$ было нормируемым, необходимо и достаточно, qтобы оно было отделимым и обладало выпуклой ограниченной окрестностью нуля (теорем а Колмогоров а).



\begin{enumerate}

  \setcounter{enumi}{4}

  \item Пусть $n$ - натуральное число, $I_{n}$ - множество, содержащее $n$ элементов и $K_{0}^{n}=K_{0}^{I_{n}}$. Топология в $K_{0}^{n}$ определяется нормой $\|x\|=\sum_{i=1}^{n}\left|x_{i}\right|$, где символ $|\cdot|$ обозначает норму в $K$. Если поле $K$ полно, то всякое $n$ мерное Т. в. п. над $K$ изоморфно пространству $K_{0}^{n}$ (при $n=1$ это верно и без предположения полноты $K$ ). Если поле $K$ локально компактно, то для того, чтобы отделимое Т. в. п. над $K$ было конечномерным, необходимо и достаточно, чтобы оно обладало предкомпактной окрестностью нуля (т е о р е м а Т и х о о о в a).

\end{enumerate}



Т. в. п. наз. м е т р и 3 у е м ы, если его топология может быть задана нек-рой метрикой (среди таких метрик всегда существует метрика, инвариантная относительно сдвигов). Для метризуемости Т. в. І. необходимо и достаточно, чтобы оно было отделимым и обладало счетной базой окрестностей нуля.



\begin{enumerate}

  \setcounter{enumi}{5}

  \item Пусть $(E, \tau)$ - Т. в. п., $E_{1}$ - векторное подпространство векторного пространства $E, \tau_{1}$ - топология, индуцированная в $E_{1}$ топологией $E$. Топология $\tau_{1}$ согласуется со структурой векторного пространства $E_{1}$; Т. в. п. $\left(E_{1}, \tau_{1}\right)$ наз. т о по л о г и ч е с ки м ве кто рным под п рост р а нс твом Т. в. п. $(E, \tau)$. Если $u-$ база (соответственно, предбаза) окрестностей нуля в $(E, \tau)$, то множество $\left\{V \cap E_{1}: V \in \mathcal{U}\right\}$ образует базу (соответственно, предбазу) окрестностей нуля в $\left(E_{1}, \tau_{1}\right)$. Если $(E, \tau)$ отделимо (соответственно, метризуемо, локально выпукло), то и $\left(E_{1}, \tau_{1}\right)$ таково же. Если топология $\tau$ задается нек-рым множеством полунорм, то топология $\tau_{1}$ определяется сужениями этих полунорм на $E_{1}$.



  \item Пусть $(E, \tau)$ и $E_{1}-$ те же, что в предыдущем пункте, и $E \mid E_{1}$ - факторпространство векторного пространства $E$ по его подпространству $E_{1}$. Фактортошология $\tau_{2}$ в $E \mid E_{1}$ согласуется со структурой векторного пространства в $E \mid E_{1} ;$ Т. в. П. $\left(E \mid E_{1}, \tau_{2}\right)$ наз. т о П о л о г ическим в е к т орным фа к тор п р о с т р анс т в о м Т. в. П. $(E, \tau)$ по $E_{1}$. (По определению фак- тортопологии, множество $V \subset E \mid E_{1}$ открыто в $\tau_{2}$ в точности тогда, когда его прообраз относительно канонич. отображения $E \rightarrow E \mid E_{1}$ открыт в $(E, \stackrel{\circ}{\tau})$. Если $U-$ база окрестностей нуля в $E$, то множество образов ее әлементов относительно канонич. отображения $E \rightarrow E \mid E_{1}$ образуют базу окрестностей нуля в $\left(E \mid E_{1}, \tau_{2}\right)$ (для предбаз это, вообще говоря, не так). Для того чтобы Т. в. п. $\left(E \mid E_{1}, \tau_{2}\right)$ было отделимым, необходимо и достаточно, чтобы подпространство $E_{1}$ было замкнуто в $(E, \tau)$. Если $\{\overline{0}\}$-замыкание одноточечного множества $\{0\}$ в $(E, \tau)$, то (отделимое) топологпч. векторное факторпространство $E /\{\overline{0}\}$ наз. о т д е ли м ы м Т. в. п., а с с о ц и и р о в а ны ы м с Т. в. п. $E ;$ конечно, если само $E$ отделимо, то ассоциированное с ним отделимое Т. в. п. ему изоморфно. Если $E$ локально выпукло (соответственно, если $E$ метризуемо, а $E_{1}$ замкнуто; если $E$ метризуемо и полно), то $E \mid E_{1}$ локально выпукло (соответственно, метризуемо; полно); однако $E$ может быть полным (неметризуемым) без того, чтобы (даже отделимое и метризуемое) его топологическое векторное факторпространство было полно (см. ниже).



  \item Пусть $\mathscr{F}$ - векторное пространство всех измеримых по Лебегу действительных функций на $[0,1], \mu_{l}$ мера Лебега на этом отрезке и для каждого $n \in \mathbb{Z}+$



\end{enumerate}



\[

V_{n}=\left\{f \in \mathscr{F}, \mu_{l}\left\{t \in[0,1]|f(t)|>\frac{1}{n+1}\right\}<\frac{1}{n+1}\right\} .

\]



Множество $\quad l=\left\{V_{n}: n \in N\right\}$ образует базис фильтра в $\mathscr{F}$, обладающий (определенными выше) свойствами (1) и (2); пусть $\tau-$ согласующаяся со структурой векторного пространства топология в $\mathscr{F}$, для к-рой $\mathcal{U}$ является базой окрестностей нуля, и Ғ - ассоциированное с $(\mathscr{T}, \tau)$ отделимое Т. в. П. (само $(\mathscr{F}, \tau)$ неотде-



\includegraphics[max width=\textwidth, center]{2023_07_09_4c135b59b085122f7e8fg-1}

ло; его можно отождествить - как векторное пространство - с пространством классов $\mu_{l}$-эквивалентных $\mu_{l}$-измеримых действительных функций на $[0,1]$; сходимость последовательностей в пространстве $(\mathscr{F}, \tau)$ (соответственно, в пространстве $\mathscr{F}_{0}$ ) совпадает со сходимостью по мере (в первом случае - индивидуальных функций, а во втором - классов $\mu_{l}$-эквивалентности таких функций).



Всюду далее предполагается, что $K=\mathbb{R}$ или $K=\mathbb{C}$.



\begin{enumerate}

  \setcounter{enumi}{8}

  \item Пусть $S=S\left(\mathbb{R}^{n}\right)$ - векторное пространство всех бесконечно диффференцируемых функций $\varphi$, определенных на $\mathbb{R}^{n}$, принимающих значения в $K$ и удовлетворяющих следующему условию [далее $\left.t=\left(t_{1}, \ldots, t_{n}\right) \in \mathbb{R}^{n}\right]$ : для всех $k, r \in \mathbb{Z}_{+}$

\end{enumerate}



\[

p_{r k}(\varphi) \equiv \max \left(1+\|t\|^{r}\right)\left\|\varphi^{(k)}(t)\right\|<\infty,

\]



где



\[

\|t\|=\left(\sum\left|t_{i}\right|^{2}\right)^{1 / 2}

\]



\[

\left\|\varphi^{(k)}(t)\right\|=\max \left\{\left|\frac{\partial^{k} \varphi(t)}{\partial t_{1}^{k_{1}} \ldots \partial_{t_{n}}^{k_{n}}}\right|: k_{1}+\ldots+k_{n}=k\right\} .

\]



Наделенное топологией $\tau_{S}$, задаваемой семейством норм $p_{r k}$, определяемых предыдущим равенством, $S$ становится полным метризуемым локально выпуклым пространством (такие пространстваназ. ІІ р о с т р а н с тв а м и Ф р е ше). Пространство $\left(S_{i}, \tau_{S}\right)$ играет важную роль в теории обобщенных функций. Интересно, что на $S$ не существует никакой нормы, превращающей $S$ в банахово пространство, в топологии к-рого каждая из функций $\varphi \mapsto \varphi(t), \quad S \rightarrow K(t \in \mathbb{R})$ непрерывна (в частности, л. в. П. $\left(S, \tau_{S}\right)$ ненормируемо).



Некоторые методы построения Т. в. п. 1. П р о е кти в ны е топо ло г и и. Пусть $E$ - векторное пространство и для каждого $\alpha$ из нек-рого множества $\mathfrak{A}$ индексов $g_{\alpha}$ - линейное отображение $E$ в Т. в. п. $E_{\alpha}$; тогда в $E$ среди всех топологиї, для к-рых непрерывны все отображения $g_{\alpha}$, существует самая слабая





\end{document}